% ==============================================================================
% Paper 87 (EN): The Bateman-Horn Conjecture: Prime Values of Polynomials
% Author: Michael Fuhrmann
% Project: Specialist Mathematics Project, Build 167
% Date: 2026-03-12
% ==============================================================================

\documentclass[12pt,a4paper]{amsart}

\usepackage{amsmath,amssymb,amsthm,mathtools,enumitem,booktabs,array,hyperref}
\usepackage[utf8]{inputenc}
\usepackage[T1]{fontenc}
\usepackage{geometry}
\geometry{margin=2.5cm}

% --------------------------------------------------------------------------
% Theorem environments
% --------------------------------------------------------------------------
\newtheorem{theorem}{Theorem}[section]
\newtheorem{lemma}[theorem]{Lemma}
\newtheorem{corollary}[theorem]{Corollary}
\newtheorem{proposition}[theorem]{Proposition}
\theoremstyle{definition}
\newtheorem{definition}[theorem]{Definition}
\newtheorem{example}[theorem]{Example}
\theoremstyle{remark}
\newtheorem{remark}[theorem]{Remark}
\newtheorem{conjecture}[theorem]{Conjecture}

% --------------------------------------------------------------------------
% Macros
% --------------------------------------------------------------------------
\newcommand{\Z}{\mathbb{Z}}
\newcommand{\Q}{\mathbb{Q}}
\newcommand{\R}{\mathbb{R}}
\newcommand{\C}{\mathbb{C}}
\newcommand{\F}{\mathbb{F}}
\newcommand{\N}{\mathbb{N}}
\newcommand{\li}{\mathrm{li}}
\DeclareMathOperator{\ord}{ord}
\newcommand{\BH}{\mathfrak{S}}

\hypersetup{
  colorlinks=true,
  linkcolor=blue,
  citecolor=blue,
  urlcolor=blue,
  pdftitle={The Bateman-Horn Conjecture: Prime Values of Polynomials},
  pdfauthor={Michael Fuhrmann}
}

% ==============================================================================
\begin{document}
% ==============================================================================

\title{The Bateman-Horn Conjecture:\\
  Prime Values of Polynomials}

\author{Michael Fuhrmann}
\address{Specialist Mathematics Project, Build 167}
\email{specialist-maths@project.local}
\date{2026-03-12}

\begin{abstract}
The Bateman-Horn conjecture, formulated in 1962, provides a quantitative prediction
for the density of integers $n \leq x$ for which $f_1(n), f_2(n), \ldots, f_k(n)$ are
simultaneously prime, where $f_1,\ldots,f_k \in \Z[x]$ are distinct irreducible
polynomials satisfying an obvious necessary condition (the \emph{product condition}).
The conjectured density is
\[
  \#\{n \leq x : f_i(n) \text{ prime for all } i\} \;\sim\;
  \BH(f_1,\ldots,f_k) \cdot \frac{x}{(\log x)^k},
\]
where $\BH(f_1,\ldots,f_k)$ is the \emph{singular series}
(an Euler product over all primes).
This unifies many classical conjectures: the Bunyakovsky conjecture ($k=1$, one quadratic),
the twin prime conjecture ($f_1 = n, f_2 = n+2$), Sophie Germain primes,
Dickson's conjecture, and Schinzel's Hypothesis H.
Proven cases include Dirichlet's theorem on primes in arithmetic progressions
(linear $f_i$ in a progression), Green-Tao's theorem (primes in arithmetic progressions),
and Maynard-Tao's theorem (bounded prime gaps).
The general Bateman-Horn conjecture remains open.
\end{abstract}

\maketitle

\tableofcontents

% ==============================================================================
\section{Introduction}
% ==============================================================================

A fundamental question in number theory is: for which polynomials $f \in \Z[x]$ are
there infinitely many $n \in \N$ such that $f(n)$ is prime?
More generally, given a system $f_1, \ldots, f_k \in \Z[x]$, when are there infinitely
many $n$ such that all $f_i(n)$ are simultaneously prime, and how often does this occur?

The Bateman-Horn conjecture \cite{BH1962} gives a unified quantitative answer:
under a natural local (``product'') condition, such $n$ exist and their density is
predicted by an Euler product (the singular series).

\subsection{Historical background}

The problem of prime values of polynomials has a rich history:
\begin{itemize}
  \item \textbf{Dirichlet (1837)}: Proved that $f(n) = an + b$ ($\gcd(a,b)=1$) represents
    infinitely many primes. This is the only family of linear polynomials fully resolved.
  \item \textbf{Bunyakovsky (1857)}: Conjectured that any irreducible $f \in \Z[x]$ with
    positive leading coefficient and no fixed prime divisor represents infinitely many primes.
    Open for $\deg f \geq 2$.
  \item \textbf{Hardy-Littlewood (1923)} \cite{HL1923}: Conjectured quantitative density
    for prime pairs $(p, p+2k)$ (prime constellations).
  \item \textbf{Dickson (1904)} \cite{Dickson1904}: Conjectured that any admissible
    system of linear forms represents infinitely many prime tuples.
  \item \textbf{Schinzel-Sierpiński (1958)} \cite{Schinzel1958}: Hypothesis H unifies
    all polynomial conjectures.
  \item \textbf{Bateman-Horn (1962)} \cite{BH1962}: Quantitative version of Hypothesis H.
\end{itemize}

\subsection{Structure of this paper}
Section~\ref{sec:setup} introduces the basic setup and definitions.
Section~\ref{sec:singular} defines the singular series.
Section~\ref{sec:special} covers special cases: Bunyakovsky, twin primes, Sophie Germain.
Section~\ref{sec:schinzel} discusses Schinzel's Hypothesis H.
Section~\ref{sec:hardy_lw} reviews Hardy-Littlewood conjectures.
Section~\ref{sec:dirichlet} discusses Dirichlet's theorem (the proven linear case).
Section~\ref{sec:sieve} covers sieve upper bounds.
Section~\ref{sec:green_tao} describes Green-Tao's theorem.
Section~\ref{sec:maynard} covers Maynard-Tao bounded gaps.
Section~\ref{sec:openproblems} lists open problems.
Section~\ref{sec:conclusion} concludes.

% ==============================================================================
\section{Setup and Definitions}\label{sec:setup}
% ==============================================================================

\begin{definition}[Admissible system]
  A system of distinct non-constant polynomials $f_1, \ldots, f_k \in \Z[x]$
  (with positive leading coefficients) is called \emph{admissible} if for every prime $p$,
  there exists an integer $n_0$ such that $p \nmid f_1(n_0) f_2(n_0) \cdots f_k(n_0)$.

  Equivalently: the product $\prod_{i=1}^k f_i$ does not vanish identically modulo $p$
  for any prime $p$.
\end{definition}

\begin{remark}
  The admissibility condition is clearly necessary for all $f_i(n)$ to be simultaneously
  prime for infinitely many $n$: if $p \mid f_i(n)$ for all $n$ and all $n$, then
  $f_i(n) = p$ at most finitely often.
\end{remark}

\begin{example}
  \begin{enumerate}[label=(\roman*)]
    \item $\{n, n+2\}$ is admissible: for $p=2$, one of $n, n+2$ is odd; for $p\geq 3$,
      not both $n \equiv 0$ and $n \equiv -2 \pmod{p}$.
    \item $\{n, n+1\}$ is NOT admissible: one of $n$, $n+1$ is always even.
    \item $\{n, 2n+1\}$ is admissible.
    \item $\{n^2+1\}$ is admissible (since $n^2+1 \not\equiv 0 \pmod{p}$ for all $n$
      for $p \equiv 3 \pmod{4}$; and for $p \equiv 1 \pmod{4}$, only half the residues
      give $n^2 \equiv -1$).
  \end{enumerate}
\end{example}

% ==============================================================================
\section{The Singular Series}\label{sec:singular}
% ==============================================================================

\begin{definition}[Singular series]
  For an admissible system $\{f_1,\ldots,f_k\}$ of irreducible polynomials over $\Z$,
  define the \emph{singular series} (also called the \emph{Hardy-Littlewood constant})
  \[
    \BH(f_1,\ldots,f_k) \;:=\;
    \prod_p \left(1 - \frac{1}{p}\right)^{-k} \left(1 - \frac{\omega(p)}{p}\right),
  \]
  where $\omega(p)$ is the number of solutions to
  $f_1(n)\cdots f_k(n) \equiv 0 \pmod{p}$
  in $\Z/p\Z$, i.e.,
  \[
    \omega(p) = \#\{n \in \Z/p\Z : f_1(n)\cdots f_k(n) \equiv 0 \pmod{p}\}.
  \]
\end{definition}

\begin{proposition}
  The singular series $\BH(f_1,\ldots,f_k)$ converges absolutely to a positive
  real number, provided the system is admissible.
\end{proposition}

\begin{proof}
  Admissibility ensures $\omega(p) \leq p-1$ for all $p$ (and $\omega(p) = 0$ for
  all but finitely many $p$ in the linear case). The Euler product converges since
  $\omega(p) = k + O(1/p)$ for large $p$ (by Hensel's lemma or direct estimation),
  giving factors $1 + O(1/p^2)$.
\end{proof}

\begin{definition}[Degree of the system]
  For polynomials $f_i$ of degrees $d_i$, define $D = \sum_{i=1}^k d_i$ (total degree)
  and $d = \prod_{i=1}^k d_i$ (product of degrees).
\end{definition}

% ==============================================================================
\section{The Bateman-Horn Conjecture}\label{sec:bh_conjecture}
% ==============================================================================

\begin{conjecture}[Bateman-Horn \cite{BH1962}]\label{conj:BH}
  Let $f_1, \ldots, f_k \in \Z[x]$ be distinct irreducible polynomials (with positive
  leading coefficients $a_1, \ldots, a_k$) forming an admissible system. Denote by
  $D = \sum_{i=1}^k \deg f_i$.
  Then
  \[
    \#\{1 \leq n \leq x : f_1(n),\ldots,f_k(n) \text{ all prime}\}
    \;\sim\;
    \frac{\BH(f_1,\ldots,f_k)}{\prod_{i=1}^k a_i}
    \cdot \frac{x}{(\log x)^k}
  \]
  as $x \to \infty$.
\end{conjecture}

\begin{remark}
  The factor $\prod a_i^{-1}$ accounts for the leading coefficients; for monic
  polynomials it equals $1$. For the system $\{f_1 = n, f_2 = n+2\}$ (twin primes),
  the conjecture gives the Hardy-Littlewood prime pair conjecture.
\end{remark}

\begin{remark}
  Bateman and Horn derived this conjecture by assuming that the primes $f_i(n)$ behave
  ``independently'' except for the obvious congruence obstructions at each prime $p$,
  captured by $\omega(p)$.
\end{remark}

% ==============================================================================
\section{Special Cases}\label{sec:special}
% ==============================================================================

\subsection{Bunyakovsky conjecture ($k=1$)}

\begin{conjecture}[Bunyakovsky / Bouniakowsky, 1857 \cite{Bunyakovsky1857}]
  \label{conj:bunyakovsky}
  Let $f \in \Z[x]$ be irreducible with positive leading coefficient and $\omega(p) < p$
  for all primes $p$. Then $f(n)$ is prime for infinitely many $n$.
\end{conjecture}

\begin{remark}
  The Bateman-Horn conjecture for $k=1$ gives the asymptotic
  \[
    \#\{n \leq x : f(n) \text{ prime}\} \;\sim\;
    \frac{\BH(f)}{a_d} \cdot \frac{x}{(\log x)},
  \]
  which is stronger than Bunyakovsky's mere assertion of infinitude.
\end{remark}

\begin{example}[$n^2+1$]
  For $f(n) = n^2 + 1$, the singular series is
  \[
    \BH(n^2+1) = \prod_{p \equiv 1 \pmod{4}} \frac{p-1}{p-2}
                 \cdot \prod_{p \equiv 3 \pmod{4}} 1
                 \;=\; 1.3203236\ldots
  \]
  (the Landau-Ramanujan constant, divided by $\sqrt{2}$).
  The conjecture predicts $\sim 1.3203 \cdot x / \log x$ primes of the form $n^2+1$ up to $x^2$.
  Infinitely many primes of the form $n^2+1$ is open.
\end{example}

\subsection{Twin prime conjecture ($k=2$, $f_1=n$, $f_2=n+2$)}

\begin{conjecture}[Twin prime conjecture]\label{conj:twin}
  There are infinitely many primes $p$ such that $p+2$ is also prime.
\end{conjecture}

\begin{conjecture}[Hardy-Littlewood conjecture B, special case of Bateman-Horn]
  \[
    \#\{p \leq x : p+2 \text{ prime}\} \;\sim\;
    2C_2 \frac{x}{(\log x)^2},
  \]
  where $C_2 = \prod_{p \geq 3} \frac{p(p-2)}{(p-1)^2} \approx 0.6601618\ldots$
  is the \emph{twin prime constant}.
\end{conjecture}

\begin{remark}
  The factor $2$ arises because the admissible pair $\{n, n+2\}$ has $\omega(p) = 2$
  for all odd primes $p$ (since if $p | n(n+2)$, then either $p|n$ or $p|n+2$, and
  for $p > 2$ these are distinct residues); so
  $\BH(\{n,n+2\}) = \prod_{p\geq 3} \frac{p-1}{p} \cdot \frac{p-2}{p-1} = 2C_2$.
\end{remark}

\subsection{Sophie Germain primes ($f_1 = n$, $f_2 = 2n+1$)}

\begin{conjecture}[Infinitude of Sophie Germain primes]\label{conj:sg}
  There are infinitely many primes $p$ such that $2p+1$ is also prime.
\end{conjecture}

The Bateman-Horn prediction for the count up to $x$:
\[
  \#\{p \leq x : 2p+1 \text{ prime}\}
  \;\sim\;
  \frac{2C_2 \cdot x}{(\log x)^2}
  \;\sim\; 2 \cdot 0.6601618 \cdot \frac{x}{(\log x)^2}
\]
(same constant as twin primes, by symmetry of the singular series).

\subsection{Primes $n^2 + n + 41$ (Euler's prime-generating formula)}

\begin{example}[Euler's polynomial]
  For $f(n) = n^2 + n + 41$, an irreducible polynomial with $\omega(p) = 0$ for
  $p \leq 41$ (since the discriminant $\Delta = 1 - 164 = -163$ has class number $1$),
  the polynomial represents primes for $n = 0, 1, 2, \ldots, 39$.
  The Bateman-Horn conjecture predicts infinitely many prime values (as $\omega(p) < p$
  for all $p$), but this is open.
\end{example}

% ==============================================================================
\section{Schinzel's Hypothesis H}\label{sec:schinzel}
% ==============================================================================

\begin{conjecture}[Schinzel's Hypothesis H \cite{Schinzel1958}]
  \label{conj:hypothesis_h}
  Let $f_1, \ldots, f_k \in \Z[x]$ be irreducible polynomials with positive leading
  coefficients such that the product $f_1 \cdots f_k$ satisfies the admissibility
  condition (no prime $p$ divides $f_1(n)\cdots f_k(n)$ for all $n \in \Z$).
  Then there exist infinitely many $n \in \N$ for which $f_1(n), \ldots, f_k(n)$ are
  all prime.
\end{conjecture}

\begin{remark}
  Schinzel's Hypothesis H is the qualitative version of the Bateman-Horn conjecture:
  Hypothesis H asserts infinitude, while Bateman-Horn predicts the asymptotic density.
  Hypothesis H implies:
  \begin{itemize}
    \item The twin prime conjecture (Conjecture~\ref{conj:twin});
    \item Infinitely many Sophie Germain primes (Conjecture~\ref{conj:sg});
    \item Infinitely many primes of the form $n^2+1$ (Bunyakovsky for $f=n^2+1$);
    \item Dickson's conjecture for prime constellations.
  \end{itemize}
\end{remark}

% ==============================================================================
\section{Hardy-Littlewood Conjectures}\label{sec:hardy_lw}
% ==============================================================================

\begin{conjecture}[Hardy-Littlewood Conjecture A \cite{HL1923}]
  \label{conj:HLA}
  The number of prime constellations of type $\{n, n+h_1, \ldots, n+h_{k-1}\}$ (for
  an admissible tuple $h_0=0, h_1, \ldots, h_{k-1}$) up to $x$ satisfies
  \[
    \pi(x; h_1,\ldots,h_{k-1}) \;\sim\; \BH_k \cdot \frac{x}{(\log x)^k},
  \]
  where $\BH_k = \BH(n, n+h_1, \ldots, n+h_{k-1})$ is the singular series.
\end{conjecture}

\begin{conjecture}[Hardy-Littlewood Conjecture B (prime pairs)]
  \[
    \#\{n \leq x : n \text{ and } n+2k \text{ both prime}\}
    \;\sim\; 2C_2 \prod_{p|k, p>2} \frac{p-1}{p-2} \cdot \frac{x}{(\log x)^2}.
  \]
\end{conjecture}

\begin{remark}
  Hardy-Littlewood Conjecture A follows from (and is essentially equivalent to) the
  Bateman-Horn conjecture for systems of linear polynomials.
\end{remark}

% ==============================================================================
\section{Dickson's Conjecture}\label{sec:dickson}
% ==============================================================================

\begin{conjecture}[Dickson, 1904 \cite{Dickson1904}]
  \label{conj:dickson}
  Let $a_1 n + b_1, \ldots, a_k n + b_k$ be $k$ linear polynomials with positive
  leading coefficients $a_i$ and $b_i \in \Z$, such that for no prime $p$ does $p$ divide
  $\prod_{i=1}^k (a_i n + b_i)$ for every $n \in \Z$.
  Then there are infinitely many $n$ for which all $a_i n + b_i$ are prime.
\end{conjecture}

\begin{remark}
  Dickson's conjecture is the linear case of Schinzel's Hypothesis H. Special cases:
  \begin{itemize}
    \item $k=1$: Dirichlet's theorem (proven, Theorem~\ref{thm:dirichlet}).
    \item $k=2$, $\{n, n+2\}$: twin prime conjecture (open).
    \item $k=2$, $\{n, 2n+1\}$: Sophie Germain primes (open).
    \item $k=2$, $\{n, n+d\}$ for any $d$: cousin/sexy primes (open for $d \neq 0$).
  \end{itemize}
\end{remark}

% ==============================================================================
\section{Dirichlet's Theorem: The Proven Linear Case}\label{sec:dirichlet}
% ==============================================================================

\begin{theorem}[Dirichlet, 1837 \cite{Dirichlet1837}]\label{thm:dirichlet}
  Let $a, b \in \Z$ with $\gcd(a,b) = 1$ and $a > 0$. Then there are infinitely many
  primes of the form $an + b$, $n \geq 1$.
  More precisely:
  \[
    \#\{p \leq x : p \equiv b \pmod{a}\} \;\sim\; \frac{1}{\phi(a)} \cdot \frac{x}{\log x},
  \]
  where $\phi(a)$ is Euler's totient function.
\end{theorem}

\begin{remark}
  Dirichlet's theorem is a special case of the Bateman-Horn conjecture for $k=1$,
  $f(n) = an+b$. In this case $\omega(p) = 0$ for $p \nmid a$, $\omega(p) = 0$ for $p|a$
  (by admissibility, since $\gcd(a,b) = 1$), and
  \[
    \BH(an+b) = \prod_{p\nmid a} \left(1-\frac{1}{p}\right)^{-1}\left(1-\frac{1}{p}\right)
              \cdot \prod_{p|a} \left(1-\frac{1}{p}\right)^{-1} \cdot 1
              = \prod_{p|a} \frac{p}{p-1}
              = \frac{a}{\phi(a)},
  \]
  giving $\sim \frac{a}{\phi(a)} \cdot \frac{1}{a} \cdot \frac{x}{\log x} = \frac{x}{\phi(a)\log x}$.
\end{remark}

The proof of Dirichlet's theorem uses Dirichlet characters $\chi\colon (\Z/a\Z)^\times \to \C^\times$
and Dirichlet $L$-functions:
\begin{equation}
  L(s, \chi) = \sum_{n=1}^\infty \frac{\chi(n)}{n^s} = \prod_p \frac{1}{1-\chi(p)p^{-s}},
\end{equation}
showing $L(1, \chi) \neq 0$ for non-principal $\chi$.

% ==============================================================================
\section{Sieve Upper Bounds}\label{sec:sieve}
% ==============================================================================

While the Bateman-Horn conjecture is unproven, sieve theory provides useful upper bounds.

\begin{theorem}[Brun's theorem, 1919 \cite{Brun1919}]
  The sum of reciprocals of twin primes converges:
  \[
    \sum_{p,p+2 \text{ both prime}} \left(\frac{1}{p} + \frac{1}{p+2}\right) < \infty.
  \]
\end{theorem}

\begin{remark}
  By contrast, the sum of reciprocals of all primes diverges (Euler). Brun's theorem
  does not prove infinitude but does not exclude it.
\end{remark}

\begin{theorem}[Chen, 1966/1973 \cite{Chen1973}]
  There are infinitely many primes $p$ such that $p+2$ is either prime or a product of
  two primes (a \emph{semiprime}).
  More generally, every even $2k$ can be written as $p + m$ where $p$ is prime and
  $m$ has at most two prime factors.
\end{theorem}

\begin{theorem}[Selberg sieve upper bound]
  For any admissible system $\{f_1,\ldots,f_k\}$ of linear polynomials,
  \[
    \#\{n \leq x : f_i(n) \text{ prime for all } i\} \;\leq\;
    C_k \cdot \frac{\BH(f_1,\ldots,f_k) \cdot x}{(\log x)^k}
  \]
  for an explicit constant $C_k > 1$. The Bateman-Horn conjecture asserts $C_k = 1 + o(1)$.
\end{theorem}

% ==============================================================================
\section{Green-Tao: Primes in Arithmetic Progressions}\label{sec:green_tao}
% ==============================================================================

One of the most celebrated theorems of the 21st century is due to Green and Tao.

\begin{theorem}[Green-Tao, 2008 \cite{GT2008}]\label{thm:green_tao}
  The prime numbers contain arithmetic progressions of arbitrary length.
  More precisely, for any $k \geq 1$, there exist infinitely many arithmetic
  progressions $\{a, a+d, a+2d, \ldots, a+(k-1)d\}$ of length $k$ with $d > 0$
  consisting entirely of prime numbers.
\end{theorem}

\begin{remark}
  The system $\{n, n+d, n+2d, \ldots, n+(k-1)d\}$ is a set of $k$ linear polynomials
  in $n$ (with $d$ fixed). The Green-Tao theorem proves \emph{infinitude} of solutions
  (not the full Bateman-Horn asymptotic). In terms of Schinzel's Hypothesis H / Dickson's
  conjecture, Green-Tao does not directly follow (since $d$ varies), but it is related.
\end{remark}

\begin{remark}
  Green and Tao's proof uses:
  \begin{enumerate}[label=(\roman*)]
    \item Szemerédi's theorem (arithmetic progressions in dense subsets of $\N$);
    \item A \emph{transference principle} showing that primes are ``pseudorandom'' enough
      for Szemerédi's theorem to apply;
    \item Green-Tao-Ziegler's generalized von Neumann theorem.
  \end{enumerate}
\end{remark}

\begin{theorem}[Quantitative Green-Tao, Green-Tao-Ziegler \cite{GTZ2012}]
  The number of $k$-term arithmetic progressions of primes $\leq x$ is
  \[
    \sim \frac{\BH_k \cdot x^2}{(\log x)^k},
  \]
  where $\BH_k = \prod_p \left(1-\frac{1}{p}\right)^{-(k-1)}\left(1-\frac{k}{p}\right)\left(1-\frac{1}{p}\right)^{-1}$
  for $p > k$ and $\BH_k$ corrected for $p \leq k$.
  This is consistent with the Bateman-Horn conjecture.
\end{theorem}

% ==============================================================================
\section{Maynard-Tao: Bounded Prime Gaps}\label{sec:maynard}
% ==============================================================================

\begin{theorem}[Zhang, 2013 \cite{Zhang2013}]
  There are infinitely many pairs of primes $(p, q)$ with $|p - q| \leq 70{,}000{,}000$.
\end{theorem}

\begin{theorem}[Maynard, 2015 \cite{Maynard2015}; Tao (independently, 2013)]
  \label{thm:maynard}
  For each $k \geq 2$, there exist $k$ distinct integers $h_1,\ldots,h_k$ such that
  \[
    \liminf_{n\to\infty} \max_{1 \leq i < j \leq k} (p_{n+j} - p_{n+i}) < \infty.
  \]
  In particular:
  \begin{enumerate}[label=(\roman*)]
    \item $\liminf_{n\to\infty} (p_{n+1} - p_n) \leq 600$ (Maynard 2015; improved to
      $246$ by the Polymath8b project \cite{Polymath2014}).
    \item For each $m \geq 1$, $\liminf_{n\to\infty}(p_{n+m} - p_n) < \infty$.
  \end{enumerate}
\end{theorem}

\begin{remark}
  Maynard-Tao's proof is via a \emph{multidimensional sieve} (GPY sieve).
  It proves that some element of a fixed admissible $k$-tuple is prime infinitely often
  (not all simultaneously prime, which would give the full twin prime conjecture).
\end{remark}

\begin{remark}
  The gap between Maynard-Tao and Bateman-Horn:
  \begin{itemize}
    \item Maynard-Tao: \emph{at least one} prime in an admissible $k$-tuple, infinitely often.
    \item Bateman-Horn: \emph{all} primes in an admissible $k$-tuple, infinitely often
      (and quantitatively).
  \end{itemize}
\end{remark}

% ==============================================================================
\section{Admissibility and Obstructions}\label{sec:admissibility}
% ==============================================================================

\begin{proposition}[Admissibility test for linear polynomials]
  The system $\{a_1 n + b_1, \ldots, a_k n + b_k\}$ is admissible iff for every prime $p$,
  the union $\{b_i/a_i \pmod{p} : i = 1,\ldots,k\}$ does not cover all of $\Z/p\Z$.
\end{proposition}

\begin{definition}[Admissible $k$-tuple]
  A tuple $(h_1,\ldots,h_k)$ with $h_i$ distinct non-negative integers is
  \emph{admissible} if for every prime $p \leq k$, not all residues
  $h_1,\ldots,h_k \pmod{p}$ are distinct (i.e., the tuple fits in fewer than $p$
  residue classes modulo $p$).
\end{definition}

\begin{example}[Smallest admissible $k$-tuples]
  \begin{center}
    \begin{tabular}{ll}
      \toprule
      $k$ & Smallest admissible $k$-tuple \\
      \midrule
      2 & $(0, 2)$ \\
      3 & $(0, 2, 6)$ \\
      4 & $(0, 2, 6, 8)$ \\
      5 & $(0, 2, 6, 8, 12)$ \\
      6 & $(0, 4, 6, 10, 12, 16)$ \\
      \bottomrule
    \end{tabular}
  \end{center}
\end{example}

% ==============================================================================
\section{Open Problems}\label{sec:openproblems}
% ==============================================================================

\begin{conjecture}[Bateman-Horn, full conjecture; Conjecture~\ref{conj:BH}]
  Open for all $k \geq 1$ with $\deg f_i \geq 2$, and for $k \geq 2$ with $\deg f_i = 1$.
  The only proven case is $k=1$, $\deg f = 1$ (Dirichlet).
\end{conjecture}

\begin{conjecture}[Twin prime conjecture; Conjecture~\ref{conj:twin}]
  Open. Best result: infinitely many prime gaps $\leq 246$ (Polymath8b \cite{Polymath2014}).
\end{conjecture}

\begin{conjecture}[Infinitude of primes $n^2+1$]
  Open. Sieve gives $\liminf_{n\to\infty} \Omega(n^2+1) \leq 2$ (Iwaniec \cite{Iwaniec1978}),
  where $\Omega$ counts prime factors with multiplicity. This does not prove primality.
\end{conjecture}

\begin{conjecture}[Bunyakovsky / Hypothesis H for $k=1$, $\deg \geq 2$]
  Open. Not a single irreducible polynomial of degree $\geq 2$ is known to represent
  infinitely many primes.
\end{conjecture}

% ==============================================================================
\section{Conclusion}\label{sec:conclusion}
% ==============================================================================

The Bateman-Horn conjecture provides a beautiful and unified quantitative framework for
the distribution of prime values of polynomial systems. Its special cases -- the twin
prime conjecture, Sophie Germain primes, Bunyakovsky's conjecture, Dickson's conjecture,
and Schinzel's Hypothesis H -- are all among the most tantalizing open problems in
number theory. The only fully proven case in the Bateman-Horn family is Dirichlet's
theorem (linear, single polynomial). Green-Tao proved infinite arithmetic progressions
of primes (consistent with Bateman-Horn for linear tuples), and Maynard-Tao established
bounded prime gaps, but neither resolves the general conjecture.

The gap between what is known (sieve upper bounds, Dirichlet, Green-Tao, Maynard-Tao)
and the full Bateman-Horn conjecture represents one of the most profound challenges
in analytic number theory.

\begin{thebibliography}{99}

\bibitem{BH1962}
P.\,T.\ Bateman, R.\,A.\ Horn,
\textit{A heuristic asymptotic formula concerning the distribution of prime numbers},
Math.\ Comp.\ \textbf{16} (1962), 363--367.

\bibitem{Brun1919}
V.\ Brun,
\textit{La série $1/5+1/7+1/11+1/13+\cdots$ où les dénominateurs sont ``nombres premiers jumeaux'' est convergente ou finie},
Bull.\ Sci.\ Math.\ \textbf{43} (1919), 100--104, 124--128.

\bibitem{Bunyakovsky1857}
V.\ Bunyakovsky,
\textit{Nouveaux théorèmes relatifs à la distinction des nombres premiers et à la
décomposition des entiers en facteurs},
Mém.\ Acad.\ Sci.\ St.\ Pétersbourg \textbf{6} (1857), 305--329.

\bibitem{Chen1973}
J.\,R.\ Chen,
\textit{On the representation of a larger even integer as the sum of a prime and the
product of at most two primes},
Sci.\ Sinica \textbf{16} (1973), 157--176.

\bibitem{Dickson1904}
L.\,E.\ Dickson,
\textit{A new extension of Dirichlet's theorem on prime numbers},
Messenger Math.\ \textbf{33} (1904), 155--161.

\bibitem{Dirichlet1837}
P.\,G.\,L.\ Dirichlet,
\textit{Beweis des Satzes, daß jede unbegrenzte arithmetische Progression, deren erstes
Glied und Differenz ganze Zahlen ohne gemeinschaftlichen Factor sind, unendlich viele
Primzahlen enthält},
Abh.\ Königl.\ Preuss.\ Akad.\ Wiss.\ (1837), 45--81.

\bibitem{GT2008}
B.\ Green, T.\ Tao,
\textit{The primes contain arbitrarily long arithmetic progressions},
Ann.\ of Math.\ (2) \textbf{167} (2008), no.\ 2, 481--547.

\bibitem{GTZ2012}
B.\ Green, T.\ Tao, T.\ Ziegler,
\textit{An inverse theorem for the Gowers $U^{s+1}[N]$-norm},
Ann.\ of Math.\ (2) \textbf{176} (2012), no.\ 2, 1231--1372.

\bibitem{HL1923}
G.\,H.\ Hardy, J.\,E.\ Littlewood,
\textit{Some problems of ``Partitio Numerorum'' III: On the expression of a number as a
sum of primes},
Acta Math.\ \textbf{44} (1923), 1--70.

\bibitem{Iwaniec1978}
H.\ Iwaniec,
\textit{Almost-primes represented by quadratic polynomials},
Invent.\ Math.\ \textbf{47} (1978), no.\ 2, 171--188.

\bibitem{Maynard2015}
J.\ Maynard,
\textit{Small gaps between primes},
Ann.\ of Math.\ (2) \textbf{181} (2015), no.\ 1, 383--413.

\bibitem{Polymath2014}
D.\,H.\,J.\ Polymath,
\textit{Variants of the Selberg sieve, and bounded intervals containing many primes},
Res.\ Math.\ Sci.\ \textbf{1} (2014), Art.\ 12.

\bibitem{Schinzel1958}
A.\ Schinzel, W.\ Sierpiński,
\textit{Sur certaines hypothèses concernant les nombres premiers},
Acta Arith.\ \textbf{4} (1958), 185--208.

\bibitem{Zhang2013}
Y.\ Zhang,
\textit{Bounded gaps between primes},
Ann.\ of Math.\ (2) \textbf{179} (2014), no.\ 3, 1121--1174.

\end{thebibliography}

\end{document}
